\begin{abstract}
Financial multi-agent systems built on large language models improve themselves exclusively through verbal reinforcement, reflecting on past trades in natural language, so risk-adjusted outcomes never update model parameters. When parameters \emph{are} updated, the target is directional accuracy, which induces reward hacking: hit-rate rises while cumulative return falls. We introduce \textsc{FinOPD}, which closes this loop through on-policy self-distillation conditioned on hindsight risk-adjusted PnL, Shapley-weighted credit assignment across specialist agents, and a non-parametric belief store. At the decision layer, Regime-Adaptive Signal Weighting and Edge-Gated Abstention restrict trading to market states where the system holds a measurable edge. Under strict post-cutoff evaluation on 2025 Dow-30 equities, against an identical-protocol suite of real trained time-series models and multi-trade agent baselines, \textsc{FinOPD} attains the highest portfolio Sharpe ratio (1.92), cumulative return, and Calmar ratio of any method, and the most stable risk-adjusted performance across evaluation windows, while a hindsight-conditioned teacher characterizes the performance ceiling (Sharpe 2.96) toward which self-evolution drives the student.
\end{abstract}
% 中文翻译：基于大语言模型的金融多智能体系统仅通过语言强化自我改进——用自然语言反思过去的交易——风险调整后的结果从未更新模型参数。当参数确实被更新时，目标是方向准确率，这诱发 reward hacking：命中率上升而累计收益下降。我们提出 FinOPD，通过以事后风险调整 PnL 为条件的在线策略自蒸馏、跨专家 agent 的 Shapley 加权信用分配和非参数 belief 存储来闭合这个回路。在决策层，Regime-Adaptive Signal Weighting 和 Edge-Gated Abstention 将交易限制在系统持有可测 edge 的市场状态。在 2025 年 Dow-30 股票上以严格 post-cutoff 评估，对由真实训练时序模型与多笔交易 agent 基线构成的同协议套件，\textsc{FinOPD} 取得最高组合夏普（1.92）、累计收益与 Calmar，且跨评测窗口风险调整表现最稳定，同时事后条件化的 teacher 刻画了自进化驱动 student 逼近的性能上界（夏普 2.96）。
